\subsection*{Two-line visual questions}
We used 16,632 clauses and 507 distinct predicate glosses from the Kyoto Classical Chinese treebank (Universal Dependencies 2.16). The released image writer, its person-reading adapter and the joint verb--object reader remained frozen. A new 150,016-parameter network projects the 160-dimensional predicate state and the pictured English question into 96-dimensional unit vectors. Two five-cell convolutions allow a question position to depend on adjacent letters. For Chinese state $i$ and English question $j$ the score is $s_{ij}=\max_{t\leq L_j}\langle u_i,v_{jt}\rangle$. With $P_i$ the batch items carrying the same registered action, the training criterion averages
\[
-\log\frac{\sum_{j\in P_i}\exp(s_{ij}/\tau)}
                {\sum_j\exp(s_{ij}/\tau)}
\]
over both retrieval directions at $\tau=0.055$. Training mixes random clauses and pairs of training instances sharing the same predicate bitmap with different glosses. The matched control permutes pictured action labels among sampled source contexts. Both fits used 1,740 AdamW steps on one RTX 4090 D. The aligned fit took 7.39 seconds and reached 79.6 MiB in allocated device tensors, excluding parent pretraining, corpus preparation and evaluation. The frozen English writer emits physical glyphs recurrently and stops on its own.

We paired eligible clauses within the same book and excluded identical action/object gloss strings across the pair. The two held frames contain the registered action name. The additional wording test was defined after literal scoring and before alternative questions were evaluated. Lean checks that, for integer comparison scores, $m$ times the number of wrong or tied comparisons is bounded by their summed hinge penalties for a positive margin $m$. This is a finite-sample property of measured scores; test accuracy and pixel fidelity require separate observations. The data release includes source identifiers, masks of eligible pairs, signed score differences and original output images. The textual interpretation of the disputed exemption passage is documented in the Academia Sinica edition of the commentary on chapter 68 of the \textit{Shiji}.\footnote{\url{https://www11.ihp.sinica.edu.tw/storage/w2_file/3806EaymdsB.pdf}}
